% !TEX program = xelatex
% 공통수학2 · Ⅰ. 도형의 방정식 · 3. 도형의 이동 중단원 마무리 + 대단원 총정리 (17차시)
% 학생용 활동지 (답안 없음)
\documentclass[11pt,a4paper]{article}
\usepackage{kotex}
\usepackage{iftex}
\ifPDFTeX\else
  \setmainhangulfont{Noto Serif CJK KR}
  \setsanshangulfont{Noto Sans CJK KR}
\fi
\usepackage[margin=2.1cm]{geometry}
\usepackage{parskip}
\usepackage{enumitem}
\usepackage{array}
\usepackage{tabularx}
\usepackage{xcolor}
\usepackage{amssymb}
\usepackage{tikz}
\usepackage[most]{tcolorbox}
\usepackage{fancyhdr}
\usepackage{titlesec}

\definecolor{goldc}{HTML}{B07C22}
\definecolor{goldstrong}{HTML}{8A5F16}
\definecolor{indigoc}{HTML}{2B3B57}
\definecolor{indigosoft}{HTML}{6E82A6}
\definecolor{linec}{HTML}{D6DACB}
\definecolor{surface2}{HTML}{E4E8DC}
\definecolor{writeline}{HTML}{C7CCBA}
\definecolor{inksoft}{HTML}{52604F}

\titleformat{\section}{\normalfont\sffamily\bfseries\large\color{indigoc}}{\thesection}{0.6em}{}
\titlespacing{\section}{0pt}{14pt}{6pt}
\newtcolorbox{goalbox}{enhanced, breakable, colback=white, colframe=goldc, boxrule=0.5pt, leftrule=3pt, arc=2pt, left=10pt, right=10pt, top=8pt, bottom=8pt}
\newtcolorbox{sheetbox}{enhanced, breakable, colback=white, colframe=linec, boxrule=0.6pt, arc=3pt, left=11pt, right=11pt, top=9pt, bottom=9pt}
\newcommand{\writespace}[1]{%
  \par\nobreak\vspace{3pt}
  \foreach \n in {1,...,#1} {%
    \noindent\textcolor{writeline}{\rule{\linewidth}{0.4pt}}\par\vspace{0.62cm}%
  }
}
\newcommand{\fillblank}[1]{\makebox[#1]{\hrulefill}}
\pagestyle{fancy}\fancyhf{}\renewcommand{\headrulewidth}{0pt}
\fancyfoot[L]{\footnotesize\color{inksoft} 공통수학2 · 2026학년도 2학기 · 지평선고등학교}
\fancyfoot[R]{\footnotesize\color{inksoft} 다음 시간 --- Ⅱ. 집합과 명제}

\begin{document}

\noindent
\begin{minipage}[b]{0.62\linewidth}
  {\sffamily\footnotesize\color{indigosoft} 공통수학2 · Ⅰ. 도형의 방정식}\\[2pt]
  {\sffamily\bfseries\Large 도형의 이동 마무리 \& 대단원 총정리 \textperiodcentered\ 17차시}
\end{minipage}\hfill
\begin{minipage}[b]{0.36\linewidth}
  \raggedleft\small\color{inksoft}
  반\ \fillblank{1.6cm} \quad 번호\ \fillblank{1.6cm}\\[4pt]
  이름\ \fillblank{3.6cm}
\end{minipage}

\vspace{4pt}
\noindent{\color{goldc}\rule{\linewidth}{2.2pt}}
\vspace{10pt}

\begin{goalbox}
\textbf{오늘의 목표}\\
도형의 이동을 정리하고, 대단원 Ⅰ 전체를 하나의 흐름으로 되짚어 본다.
\end{goalbox}

\section{대단원 지도 채우기}

\begin{sheetbox}
\begin{tabularx}{\linewidth}{@{}>{\bfseries\small}p{5cm}X@{}}
\toprule
\textbf{중단원} & \textbf{핵심 한 줄로 정리하면?} \\
\midrule
1. 평면좌표와 직선의 방정식 & \fillblank{6cm} \\[10pt]
2. 원의 방정식 & \fillblank{6cm} \\[10pt]
3. 도형의 이동 & \fillblank{6cm} \\
\bottomrule
\end{tabularx}
\end{sheetbox}

\section{도형의 이동 마무리 문제}

\begin{sheetbox}
\textbf{1.} 점 $(2,-1)$을 $x$축의 방향으로 $a$만큼, $y$축의 방향으로 $b$만큼 평행이동한 점의 좌표가 $(-5,2)$일 때, 상수 $a$, $b$의 값을 구하시오.
\writespace{2}

\textbf{2.} 다음 방정식이 나타내는 도형을 $x$축의 방향으로 $6$만큼, $y$축의 방향으로 $-2$만큼 평행이동한 도형의 방정식을 구하시오.\\
\hspace*{1.6em}⑴ $2x+3y-1=0$ \qquad ⑵ $x^2+(y-1)^2=9$
\writespace{3}

\textbf{3.} 점 $(-4,3)$을 $x$축, $y$축, 원점, 직선 $y=x$에 대하여 각각 대칭이동한 점의 좌표를 구하시오.
\writespace{3}

\textbf{4.} 원 $(x+5)^2+(y-4)^2=1$을 $x$축, $y$축, 원점, 직선 $y=x$에 대하여 각각 대칭이동한 원의 방정식을 구하시오.
\writespace{4}
\end{sheetbox}

\section{대단원 Ⅰ 평가 문제 (선택 풀이)}

\begin{sheetbox}
\textbf{5.} 직선 $3x-y-8=0$과 평행하고 점 $(3,-3)$을 지나는 직선이 점 $\left(\dfrac23, k\right)$를 지날 때, $k$의 값을 구하시오.
\writespace{3}

\textbf{6.} 방정식 $x^2+y^2+6x-4y+k+1=0$이 나타내는 도형이 원이 되도록 하는 실수 $k$의 값의 범위를 구하시오.
\writespace{3}
\end{sheetbox}

\section{대단원 성찰}

\begin{sheetbox}
\begin{minipage}[t]{0.48\linewidth}
{\footnotesize\color{inksoft} 이 대단원(Ⅰ. 도형의 방정식)에서 가장 어려웠던 순간은?}
\writespace{3}
\end{minipage}\hfill
\begin{minipage}[t]{0.48\linewidth}
{\footnotesize\color{inksoft} 가장 `아하!' 했던 순간은?}
\writespace{3}
\end{minipage}
\end{sheetbox}

\section{마무리 성찰}

\begin{sheetbox}
\begin{minipage}[t]{0.48\linewidth}
{\footnotesize\color{inksoft} 오늘 배운 것을 한 문장으로}
\writespace{2}
\end{minipage}\hfill
\begin{minipage}[t]{0.48\linewidth}
{\footnotesize\color{inksoft} 감사 일기 / 유념 한 줄}
\writespace{2}
\end{minipage}
\end{sheetbox}

\end{document}
